\scanpage{41}
1) Cas $n$ pair, $\rho=2\widetilde\rho$, $\sigma=\widetilde\sigma$, l'équation s'écrit
\[
\lambda\sigma(1+\deleted{\unsure{2}}\lambda\widetilde\rho)=0
\]
Si $\sigma$ \deleted{\unsure{est}}\inserted{et $\widetilde\rho$} inversibles,
la seule solution \uline{inversible}
(\inserted{i.e.} correspondant à une section de $U_\rho$ partout $\ne1$)
est donnée par
\[
\lambda=-\frac{1}{\deleted{\unsure{2}}\widetilde\rho}
\qquad\text{d'où}\qquad
\underbrace{u=\operatorname{id}-\frac{1}{\deleted{\unsure{2}}\unsure{\widetilde\rho}}\Pi'\otimes\Pi}
_{\uline{\text{antipodisme}}}.
\]
\deleted{Si $\uline{\sigma=0}$ $U_\rho$ invarie $f$, on a que}
D'autre part conditions équivalentes
\deleted{\unsure{on a}}
\diagram{\begin{tikzcd}[ampersand replacement=\&,column sep=0pt,row sep=small]
\& \text{b) }U_\rho\text{ invarie }f^{\mathcal E}\\
\& \text{\unsure{c)} }U_\rho\text{ invarie }Q_X^{\deleted{\unsure{\mathcal E}}}\\
\& \text{a) }\sigma=0\arrow[uu,bend left=65,start anchor=west,end anchor=west]
\end{tikzcd}}

Enfin, si $\widetilde\rho=0$, \deleted{\unsure{alors}}\inserted{ces} conditions équivalentes encore à la suivante
\deleted{$\unsure{\text{c) }\sigma=0}$}

d) $\exists$ loc. section \deleted{\unsure{différente}} de $U_\rho$ partout
$\ne1$, qui invarie $f^{\mathcal E}$.
\deleted{(\unsure{qui invarie} $Q_X$)}
\inserted{(\unsure{ou aussi} $Q_X$)}.

\deleted{\unsure{partout} $\unsure{\ne1}$}
\mspage{40}

